% This file should be rename with the speaker's last name. (e.g : smith)

% Write the author names in the order you wish them to appear
% and underline the one who will give the talk.

% Only the speaker's email address is required.

\begin{abstract_online}{Methodological issues of application of computer algebra in blended learning environment}{%
        \underline{Elena Varbanova}$^{1}$, Daniela Georgieva$^{1}$}{[elvar@yu-sofia.bg]
        }{%
        $^1$ Faculty of Applied Mathematics and Informatics, Technical University of Sofia, Sofia, Bulgaria }

We like technology for the sake of our students: it allows to transfer our knowledge and experience to them using tools and environments they are familiar with. In the application of CAS throughout the Teaching-Learning-Assessment $(TLA)$ process our main concern is to develop methodology for technology supported mathematics education [1].

It is well-known that over the centuries unique values and educational tradition have been created. We try to give contemporary/modern interpretation of the educational tradition in the country having in mind purposeful applications of CAS towards the course content, course structure, assessment model and assessment activities. The assessment activities imply the learning outcomes. Being aware of the interrelationship between the teaching, learning and assessment we design and develop teaching and learning materials based on the assessment activities. As a result we change iteratively all the three components of the TLA process. The final goal is the students to build up habits that will be later transformed into educational values.

As we teach undergraduate mathematics (subjects like Engineering Mathematics, Calculus and Numerical methods), examples of methodological approaches to selected topics will be illustrated.  The aim is to help students use prototypes, reflect on the results, understand concepts, use their imagination, work smarter not harder, master competencies [2], etc. For this purpose CAS is irreplaceable.

The ACA conferences are a kind of school for exiting and valuable collaborative work. Through any personal experience we all find out that teaching with CAS/technology is just like learning a foreign language: there is a beginning, but no end.

\textbf{Keywords} \newline{} Methodology, Computer algebra systems, Teaching-learning-assessment process

\textbf{References} \newline{}
% Model for references to books.
[1] \textsc{E. Varbanova}, About Balanced Application of CAS in Undergraduate Mathematics. In \emph{Applications of Computer Algebra, Springer Proceedings in Mathematics and Statistics}, Springer International Publishing AG 2017, 198, 2017. \newline{}
%  Model for references to articles
[2] \textsc{E. Shoikova}, \emph{Competency Based Education Development: Framework to Plan, Design and Implement Innovative CBE Programs}. Publ. House UNIBIT, Sofia, 2017. \newline{}

\end{abstract_online}

    